% !TeX root = ../main.tex
\chapter{Introduction}
\label{ch:introduction}
\section{Background and Motivation}
The purpose of this project is to develop a digital-twin-oriented training
workflow that replaces human driving demonstrations with automatically generated
supervision for a compact navigation policy. The intended outcome is a learned
LiDAR-to-steering controller that can ultimately run on a resource-constrained,
low-power development board. The central question concerns how such a controller
can be trained without requiring a person to drive and label every training
trajectory, while keeping the deployed computation small.

The intended vehicle uses a Teensy 4.0 ARM Cortex-M7 board~\cite{pjrcTeensy40}.
Training and deployment have different requirements: learning may use a host
computer, extensive simulation and information unavailable on the vehicle,
whereas onboard operation must process sensor data and produce commands within
memory, timing and energy budgets. Separating these stages permits a more capable
training system to produce a simpler deployed controller. Low-power operation
is an engineering objective of this separation; it is not an energy measurement
established by the present experiments.

In this dissertation, \emph{digital-twin-oriented training} denotes the use of
a parameterised virtual counterpart of the vehicle, sensor and control interface
to learn behaviour, generate demonstrations and test the resulting controller.
The implemented Gazebo model provides the virtual training side of this approach.
Physical calibration, live physical--virtual data exchange and experimental
confirmation of model fidelity remain future integration tasks. The present
study therefore evaluates the proposed training mechanism in simulation, without
assuming that a completed, continuously synchronised physical twin already exists.

The task is forward obstacle-avoiding cruising in static ring roads. Range-based
local control has classical precedents such as the Vector Field Histogram~\cite{borenstein1991vfh}.
Here, a learned policy converts fixed angular LiDAR bins into discrete steering
commands. The task provides a concrete setting in which to test automated
supervision and compact inference; route planning, SLAM and dynamic traffic are
outside the evaluated scope.

\section{Problem Statement}
A supervised random-forest controller needs examples of what action to take
for each sensor observation. Under a human-demonstration approach, an operator
must drive the vehicle while observations and commands are recorded. Covering
new roads, starting conditions and sensor conditions can require further
collection. Merely choosing a compact classifier does not solve this source-of-labels
problem: the fitting algorithm still needs informative, correctly paired examples.

The virtual vehicle offers a different route. Its simulated interaction provides
travel, clearance and collision feedback from which a policy can learn without
pre-existing steering labels. Once learned, that policy can drive the simulated
vehicle and supply labels automatically. Reinforcement learning (RL) is therefore
used to obtain a \emph{teacher}, rather than to perform the final random-forest
fitting directly. In the implemented workflow, Proximal Policy Optimization
(PPO) trains a neural teacher~\cite{schulman2017ppo}; supervised behaviour cloning
then fits the random forest to the teacher's LiDAR--action pairs.

This arrangement must solve three connected engineering problems. First, reward
and termination definitions must produce a teacher whose behaviour is useful
for the driving task. Second, scan and action records must preserve the
observation--decision relationship, because an incorrect temporal pairing would
create misleading supervision even for a capable teacher. Third, the student's
restricted input and model capacity must retain useful behaviour when it drives
on its own. Teacher performance does not guarantee student performance: the
student can change its future observations through its own errors~\cite{ross2011reduction}.

Replacing human demonstrations does not eliminate human engineering. Vehicle
modelling, reward design, scenario selection and acceptance criteria remain
human decisions. The completed experiments establish automated collection of
action labels; they do not quantify labour savings or establish superiority over
a human-trained controller, since no matched human-demonstration study has been
performed.

\section{Proposed Approach and Method Rationale}
The proposed solution separates learning to drive from learning a deployable
sensor-to-action mapping. A privileged teacher first learns in simulation using
map-derived distances and vehicle state. During demonstration collection, the
same simulated vehicle produces LiDAR scans and executes teacher actions. A
student then learns from these paired observations and commands, and is evaluated
in closed loop using LiDAR alone. This information separation is related to
privileged-to-sensor policy transfer~\cite{chen2019cheating}.

RL is selected because it can optimise behaviour from the consequences of
sequential actions without first requiring a labelled driving dataset. PPO is
the specific training algorithm: its clipped policy update and established
actor--critic implementation provide a practical way to train from simulated
rollouts~\cite{schulman2017ppo,raffin2021sb3}. The teacher model is a feed-forward
multilayer perceptron (MLP), with separate policy and value branches. PPO is thus
an optimisation method, MLP is the teacher's function approximator, and random
forest is the separately fitted student model. These terms describe different
parts of the workflow.

Random forests are selected as the principal student family because numerical
LiDAR bins can be partitioned by threshold tests, and tree depth and tree count
provide explicit controls over inference work~\cite{breiman2001random,sklearnForest}.
The fitted decision structures can be exported to C++ without running a neural
training framework on the vehicle. An ensemble also offers a way to combine
several partitions instead of relying on one tree. These are design motivations,
not guarantees of minimal memory, lowest power or best navigation performance.
A deeper single tree is included to test whether the ensemble's extra structure
is justified by the observed results.

The rationale is consequently a division of responsibilities: simulation supplies
repeatable interaction, RL supplies a teacher without human action labels,
behaviour cloning transfers its commands to sensor inputs, and tree inference
provides a controllable implementation for later embedded integration. Alternative
teachers and compact student families remain possible; this study evaluates one
concrete, reproducible route rather than claiming that each choice is universally
optimal.

\section{Objectives and Research Questions}
The overall objective is to demonstrate an automated path from a virtual vehicle
to a compact navigation policy, with evidence at each transition. The specific
objectives are to train and qualify an RL teacher, generate traceable demonstrations,
fit LiDAR-only students, assess the transfer of useful behaviour, and verify
exported inference before eventual board deployment.

Three questions organise the investigation:
\begin{description}
\item[RQ1: Automated supervision.] Can an RL teacher generate useful paired
LiDAR--action demonstrations in the virtual vehicle without human driving
labels, and can that collection be repeated with documented provenance?
\item[RQ2: Transfer to compact policies.] Can random-forest students retain
useful closed-loop behaviour from these demonstrations, how do they compare
with a single-tree baseline, and what limits appear under sensor loss?
\item[RQ3: Resource-constrained deployment.] Why is the selected tree
representation suitable for pursuing embedded inference, what does the verified
export establish, and which resource and physical measurements remain necessary?
\end{description}

The first question is examined through teacher qualification and two recorded
collection batches. The second uses offline prediction, primary closed-loop
comparisons and repeat runs. The third is addressed through model structure,
interface compatibility and host-side C++ agreement; actual board-level
performance remains an open experimental part of the project.

\section{Scope and Contributions}
The completed contribution is an integrated training workflow, not a new RL
algorithm or tree-learning rule. It links a parameterised virtual vehicle, a
privileged PPO teacher, automated demonstration collection, map-separated
student fitting and closed-loop evaluation. Model, map and episode records make
the supervision source traceable, and export checks connect fitted predictions
to executable C++ inference.

The primary comparison comprises 320 episodes across four controller
configurations. The single depth-six tree completes 74 of its 80 episodes,
whereas the three shallow forests complete between 62 and 67. This finding
supports evaluating the simplest useful student instead of treating a random
forest as a required winner. It also identifies a limit of the training approach:
every candidate experiences collisions under contiguous sector loss. The
comparison is configuration-specific because depth, sampling and ensemble size
are not all held constant.

Three maps supply student fitting data, while \mapname{ellipse_002} is used for
validation and selection. It is therefore not an untouched final test map~\cite{sklearnCV}.
A later batch tests consistency on the same map set without retraining. Independent-map
recovery, target-board resource consumption, physical driving and physical--virtual
fidelity have not yet been measured. The observed simulation results support the
training mechanism while leaving these broader deployment questions open.

\section{Organisation of the Dissertation}
Chapter~\ref{ch:literature} places the training strategy and compact policies in
context. Chapter~\ref{ch:system} describes the virtual counterpart and interfaces.
Chapter~\ref{ch:teacher} explains the RL choice, teacher architecture and automated
collection. Chapter~\ref{ch:student} develops the random-forest rationale, fitting
and export. Chapter~\ref{ch:results} evaluates the resulting workflow and its
limits. Chapter~\ref{ch:conclusions} answers the research questions and identifies
the remaining work. Appendix~\ref{app:evidence} records the principal evidence
sources.
